% "Como ler este relatório" -- editado à mão (fatos sobre a estrutura e os dados, não análise).
\section{Como ler este relatório}

Os capítulos 2 a 7 correspondem aos seis eixos da Política Integrada da Primeira Infância. Cada capítulo
abre com os principais achados do eixo, apresenta um indicador por seção --- com gráficos, mapas e o texto
de análise --- e termina com uma síntese. Indicadores previstos na política para os quais ainda não há
dado disponível aparecem identificados como \emph{em desenvolvimento}. As tabelas com os números completos
de cada gráfico e mapa estão nos apêndices, indicadas ao fim de cada seção. A versão interativa deste
relatório, com os mesmos indicadores e textos, está disponível no site do Instituto Pereira Passos.

Três cuidados valem para todo o relatório:

\begin{itemize}
  \item \textbf{População de referência.} Taxas do município usam as estimativas populacionais da Ripsa/Ministério
    da Saúde para o mesmo ano; taxas por bairro, região ou área programática usam a população do Censo Demográfico
    2022, a única disponível nesse nível. As duas não devem ser comparadas diretamente: o Censo 2022 tende a
    subestimar as crianças pequenas.
  \item \textbf{Faixas etárias.} Cada gráfico informa a faixa etária real do dado (por exemplo, 0 a 4 ou 0 a 5 anos),
    que nem sempre coincide com o recorte de 0 a 6 anos da política.
  \item \textbf{Proteção de dados.} Nenhuma contagem do Cadastro Único inferior a 20 crianças ou famílias é
    apresentada abaixo do nível do município; essas células aparecem vazias. Nos mapas de taxa por bairro, a
    escala de cor é limitada ao percentil 95, para que poucos bairros com números muito pequenos não escondam as
    diferenças entre os demais.
\end{itemize}
